% SPDX-License-Identifier: Apache-2.0
\section{Introduction}
\label{sec:introduction}

\IEEEPARstart{T}{wo} hardware description languages were developed
independently for the same project, each targeting the same underlying
hardware platform. The goal of this effort is a unified, consistent
language specification. This article defines the required consolidation
and resolves every open architectural and syntactic question.

Fundamentally, the two languages address complementary abstraction levels
rather than competing models. LHDL specifies system-level composition,
static configuration, and dataflow scheduling, leaving low-level protocol
timing to the compiler. TxHDL specifies cycle-accurate transaction protocols
and multi-cycle bus interactions within individual modules. Each language is
largely silent where the other provides fine detail. Four fundamental
questions require semantic decisions rather than superficial syntactic
alignment; Section~\ref{sec:conflicts} resolves each one.

\subsection{Source Material}

Seven repository files document language content, spanning initial sketches
to detailed drafts.

\begin{table*}[t]
\caption{The source material. Two languages, seven files, four spellings of
the name.}
\label{tab:sources}
\centering
\footnotesize
\begin{tabular}{@{}lrll@{}}
\toprule
Path & Lines & Language & What it is \\
\midrule
\code{spec/language.md}                  & 2241 & TxHDL & The complete TxHDL specification, 20 numbered sections \\
\code{filmil/workspace/draft-spec.md}    & 1218 & LHDL  & The long LHDL draft, prose plus EBNF \\
\code{filmil/workspace/spec.md}          &  486 & LHDL  & A condensed LHDL specification with a consolidated grammar \\
\code{filmil/workspace/lhdl-proposal.md} &  178 & LHDL  & The earliest LHDL sketch \\
\code{filmil/theses.md}                  &   94 & LHDL  & The twelve design theses the language argues from \\
\code{proposal.md}                       &   27 & both  & A prior agent's read of the repository and its next steps \\
\code{README.md}                         &    1 & none  & One sentence naming the two authors \\
\bottomrule
\end{tabular}
\end{table*}

Prior to this consolidation, the most comprehensive specification in the
repository was neither integrated into the automated build nor published.
The complete specification in \code{spec/language.md} lacked a build target,
whereas several smaller exploratory files possessed active Bazel rules.

Nomenclature also diverged across these documents. The extended LHDL draft
uses FLHDL 88 times; the condensed specification uses it 9 times; the initial
proposal uses LHdl 6 times; and \code{spec/language.md} introduces TxHDL 5
times. Because the repository, build targets, and distribution packages use
\code{txhdl}, Section~\ref{sec:decisions} adopts TxHDL as the canonical name.

\subsection{Contributions}

\begin{enumerate}
  \item An account of what each language actually specifies, and of the
        asymmetry that makes a merge possible rather than a compromise
        (Section~\ref{sec:two-languages}).
  \item Resolutions for the four semantic conflicts, of which the model of
        time is the one that shapes the language
        (Section~\ref{sec:conflicts}).
  \item A surface syntax biased toward Rust, decided construct by
        construct, which removes six defects present in the sources rather
        than restating them (Section~\ref{sec:syntax}).
  \item An inventory of every defect found by reading the two sources,
        with line numbers, including sixteen places where the LHDL grammar
        and the LHDL examples disagree
        (Section~\ref{sec:defects}).
\end{enumerate}
